\section{Questao 1 - BFS, DFS e Busca Iterativa em Profundidade}

\subsection{Formulacao do problema em espaco de estados}
\textbf{a) Definicao formal do estado}
\[
P = \langle Z, \mathcal{A}, T, s_0, G, c \rangle
\]
\begin{itemize}[leftmargin=1.5em]
  \item $Z$: conjunto de estados (salas do predio), com
  \[
  Z=\{A,B,C,D,E,F,G,H,I,J,K,L,M,N,O,P,Q,R,S\}.
  \]
  \item Cada estado $z \in Z$ representa a sala atual do agente.
  \item $\mathcal{A}$: conjunto de acoes de deslocamento entre salas adjacentes.
  \item $T$: funcao de transicao que leva da sala atual para uma sala vizinha valida.
\end{itemize}

\textbf{b) Estado inicial}
\[
s_0 = A
\]

\textbf{c) Teste de objetivo}
\[
G(z)=
\begin{cases}
\text{verdadeiro}, & \text{se } z=S \\
\text{falso}, & \text{caso contrario}
\end{cases}
\]

\textbf{d) Funcao sucessora}
\begin{itemize}[leftmargin=1.5em]
  \item A funcao sucessora devolve os vizinhos da sala atual respeitando a ordem obrigatoria do enunciado.
  \item Exemplo: $\text{Suc}(A)=[B,C,D]$, $\text{Suc}(C)=[G,H]$, $\text{Suc}(F)=[K,L]$.
\end{itemize}

\textbf{e) Funcao de custo}
\[
c(z,a,z') = 1
\]
Cada movimento entre salas adjacentes tem custo unitario. Logo, o custo total de um caminho e o numero de arestas percorridas.

\textbf{f) Representacao escolhida em Python}
\begin{lstlisting}
grafo = {
    "A": ["B", "C", "D"],
    "B": ["E", "F"],
    "C": ["G", "H"],
    "D": ["I"],
    "E": ["J"],
    "F": ["K", "L"],
    "G": ["M"],
    "H": ["N", "O"],
    "I": ["P"],
    "J": [],
    "K": ["Q"],
    "L": [],
    "M": ["R"],
    "N": [],
    "O": ["S"],
    "P": [],
    "Q": [],
    "R": [],
    "S": [],
}
estado_inicial = "A"
estado_objetivo = "S"
\end{lstlisting}

\subsection{Grafo definido no enunciado}
\begin{figure}[H]
\centering
\begin{tikzpicture}[
  >=Stealth,
  node distance=10mm,
  state/.style={draw, circle, minimum size=7mm, inner sep=0pt},
  edge/.style={->, thick}
]
  % Nivel 0
  \node[state] (A) at (0,0) {A};

  % Nivel 1
  \node[state] (B) at (2,2) {B};
  \node[state] (C) at (2,0) {C};
  \node[state] (D) at (2,-2) {D};

  % Nivel 2
  \node[state] (E) at (4,3) {E};
  \node[state] (F) at (4,1.5) {F};
  \node[state] (G) at (4,0) {G};
  \node[state] (H) at (4,-1.5) {H};
  \node[state] (I) at (4,-3) {I};

  % Nivel 3
  \node[state] (J) at (6,3.5) {J};
  \node[state] (K) at (6,2.2) {K};
  \node[state] (L) at (6,0.9) {L};
  \node[state] (M) at (6,-0.4) {M};
  \node[state] (N) at (6,-1.7) {N};
  \node[state] (O) at (6,-3.0) {O};
  \node[state] (P) at (6,-4.3) {P};

  % Nivel 4
  \node[state] (Q) at (8,2.2) {Q};
  \node[state] (R) at (8,-0.4) {R};
  \node[state] (S) at (8,-3.0) {S};

  % Arestas
  \draw[edge] (A) -- (B);
  \draw[edge] (A) -- (C);
  \draw[edge] (A) -- (D);

  \draw[edge] (B) -- (E);
  \draw[edge] (B) -- (F);

  \draw[edge] (C) -- (G);
  \draw[edge] (C) -- (H);

  \draw[edge] (D) -- (I);

  \draw[edge] (E) -- (J);
  \draw[edge] (F) -- (K);
  \draw[edge] (F) -- (L);
  \draw[edge] (G) -- (M);
  \draw[edge] (H) -- (N);
  \draw[edge] (H) -- (O);
  \draw[edge] (I) -- (P);

  \draw[edge] (K) -- (Q);
  \draw[edge] (M) -- (R);
  \draw[edge] (O) -- (S);
\end{tikzpicture}
\caption{Grafo da Q1}
\end{figure}

\subsection{Arvore de busca da BFS (com ordem de expansao)}
\begin{figure}[H]
\centering
\begin{tikzpicture}[
  >=Stealth,
  state/.style={draw, circle, minimum size=8mm, inner sep=0pt},
  edge/.style={->, thick}
]
  % Nos com numero de expansao no superscrito
  \node[state] (A) at (0,0) {$A^{1}$};

  \node[state] (B) at (2,2) {$B^{2}$};
  \node[state] (C) at (2,0) {$C^{3}$};
  \node[state] (D) at (2,-2) {$D^{4}$};

  \node[state] (E) at (4,3) {$E^{5}$};
  \node[state] (F) at (4,1.5) {$F^{6}$};
  \node[state] (G) at (4,0) {$G^{7}$};
  \node[state] (H) at (4,-1.5) {$H^{8}$};
  \node[state] (I) at (4,-3) {$I^{9}$};

  \node[state] (J) at (6,3.5) {$J^{10}$};
  \node[state] (K) at (6,2.2) {$K^{11}$};
  \node[state] (L) at (6,0.9) {$L^{12}$};
  \node[state] (M) at (6,-0.4) {$M^{13}$};
  \node[state] (N) at (6,-1.7) {$N^{14}$};
  \node[state] (O) at (6,-3.0) {$O^{15}$};
  \node[state] (P) at (6,-4.3) {$P^{16}$};

  \node[state] (Q) at (8,2.2) {$Q^{17}$};
  \node[state] (R) at (8,-0.4) {$R^{18}$};
  \node[state] (S) at (8,-3.0) {$S^{19}$};

  % Arestas da arvore BFS
  \draw[edge] (A) -- (B);
  \draw[edge, blue!70!black, very thick] (A) -- (C);
  \draw[edge] (A) -- (D);

  \draw[edge] (B) -- (E);
  \draw[edge] (B) -- (F);

  \draw[edge] (C) -- (G);
  \draw[edge, blue!70!black, very thick] (C) -- (H);

  \draw[edge] (D) -- (I);

  \draw[edge] (E) -- (J);
  \draw[edge] (F) -- (K);
  \draw[edge] (F) -- (L);
  \draw[edge] (G) -- (M);
  \draw[edge] (H) -- (N);
  \draw[edge, blue!70!black, very thick] (H) -- (O);
  \draw[edge] (I) -- (P);

  \draw[edge] (K) -- (Q);
  \draw[edge] (M) -- (R);
  \draw[edge, blue!70!black, very thick] (O) -- (S);
\end{tikzpicture}
\caption{Arvore de busca BFS com indice de expansao; caminho solucao destacado em azul}
\end{figure}

\subsection{Fronteira e visitados da BFS}
\textbf{Convencao usada:}
\begin{itemize}[leftmargin=1.5em]
  \item \textit{Fronteira} = fila de abertos (da esquerda para a direita).
  \item \textit{Visitados} = ordem de nos efetivamente expandidos.
  \item Estado inicial: fronteira $=[A]$ e visitados $=\varnothing$.
\end{itemize}

\begingroup
\footnotesize
\setlength{\tabcolsep}{4pt}
\renewcommand{\arraystretch}{1.15}
\begin{longtable}{c c p{4.9cm} p{5.9cm}}
\toprule
Passo & No expandido & Fronteira (abertos) & Visitados \\
\midrule
0 & - & [A] & [] \\
1 & A & [B, C, D] & [A] \\
2 & B & [C, D, E, F] & [A, B] \\
3 & C & [D, E, F, G, H] & [A, B, C] \\
4 & D & [E, F, G, H, I] & [A, B, C, D] \\
5 & E & [F, G, H, I, J] & [A, B, C, D, E] \\
6 & F & [G, H, I, J, K, L] & [A, B, C, D, E, F] \\
7 & G & [H, I, J, K, L, M] & [A, B, C, D, E, F, G] \\
8 & H & [I, J, K, L, M, N, O] & [A, B, C, D, E, F, G, H] \\
9 & I & [J, K, L, M, N, O, P] & [A, B, C, D, E, F, G, H, I] \\
10 & J & [K, L, M, N, O, P] & [A, B, C, D, E, F, G, H, I, J] \\
11 & K & [L, M, N, O, P, Q] & [A, B, C, D, E, F, G, H, I, J, K] \\
12 & L & [M, N, O, P, Q] & [A, B, C, D, E, F, G, H, I, J, K, L] \\
13 & M & [N, O, P, Q, R] & [A, B, C, D, E, F, G, H, I, J, K, L, M] \\
14 & N & [O, P, Q, R] & [A, B, C, D, E, F, G, H, I, J, K, L, M, N] \\
15 & O & [P, Q, R, S] & [A, B, C, D, E, F, G, H, I, J, K, L, M, N, O] \\
16 & P & [Q, R, S] & [A, B, C, D, E, F, G, H, I, J, K, L, M, N, O, P] \\
17 & Q & [R, S] & [A, B, C, D, E, F, G, H, I, J, K, L, M, N, O, P, Q] \\
18 & R & [S] & [A, B, C, D, E, F, G, H, I, J, K, L, M, N, O, P, Q, R] \\
19 & S & [] & [A, B, C, D, E, F, G, H, I, J, K, L, M, N, O, P, Q, R, S] \\
\bottomrule
\end{longtable}
\endgroup

\subsection{Resultado final da BFS}
\begin{itemize}[leftmargin=1.5em]
  \item Caminho solucao: $A \rightarrow C \rightarrow H \rightarrow O \rightarrow S$.
  \item Profundidade da solucao: 4.
  \item Custo da solucao: 4.
  \item Nos gerados: 19.
  \item Nos expandidos: 19.
\end{itemize}

\subsection{Arvore de busca da DFS (com ordem de expansao)}
\begin{figure}[H]
\centering
\begin{tikzpicture}[
  >=Stealth,
  state/.style={draw, circle, minimum size=8mm, inner sep=0pt},
  edge/.style={->, thick}
]
  \node[state] (A) at (0,0) {$A^{1}$};

  \node[state] (B) at (2,2) {$B^{2}$};
  \node[state] (C) at (2,0) {$C^{9}$};
  \node[state] (D) at (2,-2) {D};

  \node[state] (E) at (4,3) {$E^{3}$};
  \node[state] (F) at (4,1.5) {$F^{5}$};
  \node[state] (G) at (4,0) {$G^{10}$};
  \node[state] (H) at (4,-1.5) {$H^{13}$};

  \node[state] (J) at (6,3.5) {$J^{4}$};
  \node[state] (K) at (6,2.2) {$K^{6}$};
  \node[state] (L) at (6,0.9) {$L^{8}$};
  \node[state] (M) at (6,-0.4) {$M^{11}$};
  \node[state] (N) at (6,-1.7) {$N^{14}$};
  \node[state] (O) at (6,-3.0) {$O^{15}$};

  \node[state] (Q) at (8,2.2) {$Q^{7}$};
  \node[state] (R) at (8,-0.4) {$R^{12}$};
  \node[state] (S) at (8,-3.0) {$S^{16}$};

  \draw[edge] (A) -- (B);
  \draw[edge, blue!70!black, very thick] (A) -- (C);
  \draw[edge] (A) -- (D);

  \draw[edge] (B) -- (E);
  \draw[edge] (B) -- (F);
  \draw[edge] (E) -- (J);
  \draw[edge] (F) -- (K);
  \draw[edge] (F) -- (L);
  \draw[edge] (K) -- (Q);

  \draw[edge] (C) -- (G);
  \draw[edge, blue!70!black, very thick] (C) -- (H);
  \draw[edge] (G) -- (M);
  \draw[edge] (M) -- (R);
  \draw[edge] (H) -- (N);
  \draw[edge, blue!70!black, very thick] (H) -- (O);
  \draw[edge, blue!70!black, very thick] (O) -- (S);
\end{tikzpicture}
\caption{Arvore de busca DFS com indice de expansao; caminho solucao em azul}
\end{figure}

\subsection{Fronteira e visitados da DFS}
\begin{itemize}[leftmargin=1.5em]
  \item Fronteira = pilha de abertos.
  \item O topo da pilha esta na extrema direita da lista.
  \item Estado inicial: fronteira $=[A]$ e visitados $=\varnothing$.
\end{itemize}

\begingroup
\footnotesize
\setlength{\tabcolsep}{4pt}
\renewcommand{\arraystretch}{1.15}
\begin{longtable}{c c p{4.9cm} p{5.9cm}}
\toprule
Passo & No expandido & Fronteira (pilha) & Visitados \\
\midrule
0 & - & [A] & [] \\
1 & A & [D, C, B] & [A] \\
2 & B & [D, C, F, E] & [A, B] \\
3 & E & [D, C, F, J] & [A, B, E] \\
4 & J & [D, C, F] & [A, B, E, J] \\
5 & F & [D, C, L, K] & [A, B, E, J, F] \\
6 & K & [D, C, L, Q] & [A, B, E, J, F, K] \\
7 & Q & [D, C, L] & [A, B, E, J, F, K, Q] \\
8 & L & [D, C] & [A, B, E, J, F, K, Q, L] \\
9 & C & [D, H, G] & [A, B, E, J, F, K, Q, L, C] \\
10 & G & [D, H, M] & [A, B, E, J, F, K, Q, L, C, G] \\
11 & M & [D, H, R] & [A, B, E, J, F, K, Q, L, C, G, M] \\
12 & R & [D, H] & [A, B, E, J, F, K, Q, L, C, G, M, R] \\
13 & H & [D, O, N] & [A, B, E, J, F, K, Q, L, C, G, M, R, H] \\
14 & N & [D, O] & [A, B, E, J, F, K, Q, L, C, G, M, R, H, N] \\
15 & O & [D, S] & [A, B, E, J, F, K, Q, L, C, G, M, R, H, N, O] \\
16 & S & [D] & [A, B, E, J, F, K, Q, L, C, G, M, R, H, N, O, S] \\
\bottomrule
\end{longtable}
\endgroup

\subsection{Resultado final da DFS}
\begin{itemize}[leftmargin=1.5em]
  \item Caminho solucao: $A \rightarrow C \rightarrow H \rightarrow O \rightarrow S$.
  \item Profundidade da solucao: 4.
  \item Custo da solucao: 4.
  \item Nos gerados: 17.
  \item Nos expandidos: 16.
\end{itemize}

\subsection{Busca iterativa em profundidade (IDS)}
\textbf{Resumo por limite de profundidade:}
\begin{table}[H]
\centering
\begin{tabular}{c c c c}
\toprule
Limite & Objetivo encontrado? & Nos gerados & Nos expandidos \\
\midrule
0 & nao & 1 & 1 \\
1 & nao & 4 & 4 \\
2 & nao & 9 & 9 \\
3 & nao & 16 & 16 \\
4 & sim & 17 & 16 \\
\bottomrule
\end{tabular}
\caption{IDS por iteracao de limite}
\end{table}

\textbf{Ordem de visita aos nos (IDS):}
\begin{table}[H]
\centering
\begin{tabular}{c p{10.5cm}}
\toprule
Limite & Ordem de visita (expansao) \\
\midrule
0 & [A] \\
1 & [A, B, C, D] \\
2 & [A, B, E, F, C, G, H, D, I] \\
3 & [A, B, E, J, F, K, L, C, G, M, H, N, O, D, I, P] \\
4 & [A, B, E, J, F, K, Q, L, C, G, M, R, H, N, O, S] \\
\bottomrule
\end{tabular}
\caption{IDS: ordem de visita dos nos em cada iteracao}
\end{table}

\subsection{Arvores da IDS em todas as iteracoes (limites 0 a 4)}
\begin{figure}[H]
\centering
\begin{subfigure}[t]{0.32\textwidth}
\centering
\begin{forest}
for tree={
  draw,
  circle,
  minimum size=6mm,
  inner sep=0pt,
  l sep=8mm,
  s sep=6mm,
  edge={->,thick}
}
[{$A^{1}$}]
\end{forest}
\caption{Limite 0}
\end{subfigure}
\hfill
\begin{subfigure}[t]{0.32\textwidth}
\centering
\begin{forest}
for tree={
  draw,
  circle,
  minimum size=6mm,
  inner sep=0pt,
  l sep=8mm,
  s sep=6mm,
  edge={->,thick}
}
[{$A^{1}$}
  [{$B^{2}$}]
  [{$C^{3}$}, edge+={blue!70!black,very thick}]
  [{$D^{4}$}]
]
\end{forest}
\caption{Limite 1}
\end{subfigure}
\hfill
\begin{subfigure}[t]{0.32\textwidth}
\centering
\begin{forest}
for tree={
  draw,
  circle,
  minimum size=6mm,
  inner sep=0pt,
  l sep=7mm,
  s sep=4mm,
  edge={->,thick}
}
[{$A^{1}$}
  [{$B^{2}$}
    [{$E^{3}$}]
    [{$F^{4}$}]
  ]
  [{$C^{5}$}, edge+={blue!70!black,very thick}
    [{$G^{6}$}]
    [{$H^{7}$}, edge+={blue!70!black,very thick}]
  ]
  [{$D^{8}$}
    [{$I^{9}$}]
  ]
]
\end{forest}
\caption{Limite 2}
\end{subfigure}

\vspace{0.8em}

\begin{subfigure}[t]{0.49\textwidth}
\centering
\begin{forest}
for tree={
  draw,
  circle,
  minimum size=6mm,
  inner sep=0pt,
  l sep=7mm,
  s sep=4mm,
  edge={->,thick}
}
[{$A^{1}$}
  [{$B^{2}$}
    [{$E^{3}$}
      [{$J^{4}$}]
    ]
    [{$F^{5}$}
      [{$K^{6}$}]
      [{$L^{7}$}]
    ]
  ]
  [{$C^{8}$}, edge+={blue!70!black,very thick}
    [{$G^{9}$}
      [{$M^{10}$}]
    ]
    [{$H^{11}$}, edge+={blue!70!black,very thick}
      [{$N^{12}$}]
      [{$O^{13}$}, edge+={blue!70!black,very thick}]
    ]
  ]
  [{$D^{14}$}
    [{$I^{15}$}
      [{$P^{16}$}]
    ]
  ]
]
\end{forest}
\caption{Limite 3}
\end{subfigure}
\hfill
\begin{subfigure}[t]{0.49\textwidth}
\centering
\begin{forest}
for tree={
  draw,
  circle,
  minimum size=6mm,
  inner sep=0pt,
  l sep=7mm,
  s sep=4mm,
  edge={->,thick}
}
[{$A^{1}$}
  [{$B^{2}$}
    [{$E^{3}$}
      [{$J^{4}$}]
    ]
    [{$F^{5}$}
      [{$K^{6}$}
        [{$Q^{7}$}]
      ]
      [{$L^{8}$}]
    ]
  ]
  [{$C^{9}$}, edge+={blue!70!black,very thick}
    [{$G^{10}$}
      [{$M^{11}$}
        [{$R^{12}$}]
      ]
    ]
    [{$H^{13}$}, edge+={blue!70!black,very thick}
      [{$N^{14}$}]
      [{$O^{15}$}, edge+={blue!70!black,very thick}
        [{$S^{16}$}, edge+={blue!70!black,very thick}]
      ]
    ]
  ]
  [D]
]
\end{forest}
\caption{Limite 4}
\end{subfigure}
\caption{IDS: uma arvore por iteracao de limite; caminho solucao em azul}
\end{figure}

\subsection{Fronteira e visitados da IDS (iteracao limite 4)}
\begin{itemize}[leftmargin=1.5em]
  \item Nesta tabela, mostramos a iteracao que encontrou o objetivo.
  \item A fronteira e tratada como pilha (topo na extrema direita).
\end{itemize}

\begingroup
\footnotesize
\setlength{\tabcolsep}{4pt}
\renewcommand{\arraystretch}{1.15}
\begin{longtable}{c c p{4.9cm} p{5.9cm}}
\toprule
Passo & No expandido & Fronteira (pilha) & Visitados \\
\midrule
0 & - & [A] & [] \\
1 & A & [D, C, B] & [A] \\
2 & B & [D, C, F, E] & [A, B] \\
3 & E & [D, C, F, J] & [A, B, E] \\
4 & J & [D, C, F] & [A, B, E, J] \\
5 & F & [D, C, L, K] & [A, B, E, J, F] \\
6 & K & [D, C, L, Q] & [A, B, E, J, F, K] \\
7 & Q & [D, C, L] & [A, B, E, J, F, K, Q] \\
8 & L & [D, C] & [A, B, E, J, F, K, Q, L] \\
9 & C & [D, H, G] & [A, B, E, J, F, K, Q, L, C] \\
10 & G & [D, H, M] & [A, B, E, J, F, K, Q, L, C, G] \\
11 & M & [D, H, R] & [A, B, E, J, F, K, Q, L, C, G, M] \\
12 & R & [D, H] & [A, B, E, J, F, K, Q, L, C, G, M, R] \\
13 & H & [D, O, N] & [A, B, E, J, F, K, Q, L, C, G, M, R, H] \\
14 & N & [D, O] & [A, B, E, J, F, K, Q, L, C, G, M, R, H, N] \\
15 & O & [D, S] & [A, B, E, J, F, K, Q, L, C, G, M, R, H, N, O] \\
16 & S & [D] & [A, B, E, J, F, K, Q, L, C, G, M, R, H, N, O, S] \\
\bottomrule
\end{longtable}
\endgroup

\subsection{Resultado final da IDS}
\begin{itemize}[leftmargin=1.5em]
  \item Caminho solucao: $A \rightarrow C \rightarrow H \rightarrow O \rightarrow S$.
  \item Profundidade da solucao: 4.
  \item Custo da solucao: 4.
  \item Nos gerados totais (somando limites 0..4): 47.
  \item Nos expandidos totais (somando limites 0..4): 46.
\end{itemize}

\subsection{Experimento adicional: alteracao da ordem de sucessores}
\textbf{Alteracoes aplicadas (exatamente dois nos):}
\begin{itemize}[leftmargin=1.5em]
  \item $C: [G,H] \rightarrow [H,G]$
  \item $F: [K,L] \rightarrow [L,K]$
\end{itemize}

\textbf{Resultados apos reexecucao dos tres algoritmos:}
\begin{table}[H]
\centering
\small
\setlength{\tabcolsep}{4.2pt}
\renewcommand{\arraystretch}{1.15}
\begin{tabular}{>{\raggedright\arraybackslash}p{2.2cm} >{\raggedright\arraybackslash}p{4.9cm} c c c c}
\toprule
Algoritmo & Caminho & Prof. & Custo & Gerados & Expandidos \\
\midrule
BFS (original) & $A \rightarrow C \rightarrow H \rightarrow O \rightarrow S$ & 4 & 4 & 19 & 19 \\
BFS (modificado) & $A \rightarrow C \rightarrow H \rightarrow O \rightarrow S$ & 4 & 4 & 19 & 18 \\
\midrule
DFS (original) & $A \rightarrow C \rightarrow H \rightarrow O \rightarrow S$ & 4 & 4 & 17 & 16 \\
DFS (modificado) & $A \rightarrow C \rightarrow H \rightarrow O \rightarrow S$ & 4 & 4 & 15 & 13 \\
\midrule
IDS (original, total) & $A \rightarrow C \rightarrow H \rightarrow O \rightarrow S$ & 4 & 4 & 47 & 46 \\
IDS (modificado, total) & $A \rightarrow C \rightarrow H \rightarrow O \rightarrow S$ & 4 & 4 & 45 & 43 \\
\bottomrule
\end{tabular}
\caption{Comparacao numerica antes/depois da alteracao}
\end{table}

\subsection{Redesenho completo apos a alteracao de sucessores}
\textbf{Grafo apos alteracao em dois nos:}
\begin{itemize}[leftmargin=1.5em]
  \item $C: [H,G]$
  \item $F: [L,K]$
\end{itemize}

\begin{figure}[H]
\centering
\begin{tikzpicture}[
  >=Stealth,
  state/.style={draw, circle, minimum size=7mm, inner sep=0pt},
  edge/.style={->, thick}
]
  \node[state] (A) at (0,0) {A};
  \node[state] (B) at (2,2) {B};
  \node[state] (C) at (2,0) {C};
  \node[state] (D) at (2,-2) {D};

  \node[state] (E) at (4,3) {E};
  \node[state] (F) at (4,1.5) {F};
  \node[state] (G) at (4,0) {G};
  \node[state] (H) at (4,-1.5) {H};
  \node[state] (I) at (4,-3) {I};

  \node[state] (J) at (6,3.5) {J};
  \node[state] (K) at (6,2.2) {K};
  \node[state] (L) at (6,0.9) {L};
  \node[state] (M) at (6,-0.4) {M};
  \node[state] (N) at (6,-1.7) {N};
  \node[state] (O) at (6,-3.0) {O};
  \node[state] (P) at (6,-4.3) {P};

  \node[state] (Q) at (8,2.2) {Q};
  \node[state] (R) at (8,-0.4) {R};
  \node[state] (S) at (8,-3.0) {S};

  \draw[edge] (A) -- (B);
  \draw[edge] (A) -- (C);
  \draw[edge] (A) -- (D);
  \draw[edge] (B) -- (E);
  \draw[edge] (B) -- (F);
  \draw[edge] (C) -- (H);
  \draw[edge] (C) -- (G);
  \draw[edge] (D) -- (I);
  \draw[edge] (E) -- (J);
  \draw[edge] (F) -- (L);
  \draw[edge] (F) -- (K);
  \draw[edge] (G) -- (M);
  \draw[edge] (H) -- (N);
  \draw[edge] (H) -- (O);
  \draw[edge] (I) -- (P);
  \draw[edge] (K) -- (Q);
  \draw[edge] (M) -- (R);
  \draw[edge] (O) -- (S);
\end{tikzpicture}
\caption{Grafo apos alteracao da ordem de sucessores (estrutura igual, ordem local alterada em C e F)}
\end{figure}

\textbf{Arvore BFS (ordem modificada):}
\begin{figure}[H]
\centering
\begin{tikzpicture}[
  >=Stealth,
  state/.style={draw, circle, minimum size=8mm, inner sep=0pt},
  edge/.style={->, thick}
]
  \node[state] (A) at (0,0) {$A^{1}$};
  \node[state] (B) at (2,2) {$B^{2}$};
  \node[state] (C) at (2,0) {$C^{3}$};
  \node[state] (D) at (2,-2) {$D^{4}$};

  \node[state] (E) at (4,3) {$E^{5}$};
  \node[state] (F) at (4,1.5) {$F^{6}$};
  \node[state] (H) at (4,0) {$H^{7}$};
  \node[state] (G) at (4,-1.5) {$G^{8}$};
  \node[state] (I) at (4,-3) {$I^{9}$};

  \node[state] (J) at (6,3.5) {$J^{10}$};
  \node[state] (L) at (6,2.2) {$L^{11}$};
  \node[state] (K) at (6,0.9) {$K^{12}$};
  \node[state] (N) at (6,-0.4) {$N^{13}$};
  \node[state] (O) at (6,-1.7) {$O^{14}$};
  \node[state] (M) at (6,-3.0) {$M^{15}$};
  \node[state] (P) at (6,-4.3) {$P^{16}$};

  \node[state] (Q) at (8,0.9) {$Q^{17}$};
  \node[state] (S) at (8,-1.7) {$S^{18}$};
  \node[state] (R) at (8,-3.0) {R};

  \draw[edge] (A) -- (B);
  \draw[edge, blue!70!black, very thick] (A) -- (C);
  \draw[edge] (A) -- (D);
  \draw[edge] (B) -- (E);
  \draw[edge] (B) -- (F);
  \draw[edge, blue!70!black, very thick] (C) -- (H);
  \draw[edge] (C) -- (G);
  \draw[edge] (D) -- (I);
  \draw[edge] (E) -- (J);
  \draw[edge] (F) -- (L);
  \draw[edge] (F) -- (K);
  \draw[edge] (H) -- (N);
  \draw[edge, blue!70!black, very thick] (H) -- (O);
  \draw[edge] (G) -- (M);
  \draw[edge] (I) -- (P);
  \draw[edge] (K) -- (Q);
  \draw[edge] (M) -- (R);
  \draw[edge, blue!70!black, very thick] (O) -- (S);
\end{tikzpicture}
\caption{BFS apos alteracao; R permanece na fronteira quando S e encontrado}
\end{figure}

\textbf{Arvore DFS (ordem modificada):}
\begin{figure}[H]
\centering
\begin{tikzpicture}[
  >=Stealth,
  state/.style={draw, circle, minimum size=8mm, inner sep=0pt},
  edge/.style={->, thick}
]
  \node[state] (A) at (0,0) {$A^{1}$};
  \node[state] (B) at (2,2) {$B^{2}$};
  \node[state] (C) at (2,0) {$C^{9}$};
  \node[state] (D) at (2,-2) {D};

  \node[state] (E) at (4,3) {$E^{3}$};
  \node[state] (F) at (4,1.5) {$F^{5}$};
  \node[state] (H) at (4,0) {$H^{10}$};
  \node[state] (G) at (4,-1.5) {G};

  \node[state] (J) at (6,3.5) {$J^{4}$};
  \node[state] (L) at (6,2.2) {$L^{6}$};
  \node[state] (K) at (6,0.9) {$K^{7}$};
  \node[state] (N) at (6,-0.4) {$N^{11}$};
  \node[state] (O) at (6,-1.7) {$O^{12}$};

  \node[state] (Q) at (8,0.9) {$Q^{8}$};
  \node[state] (S) at (8,-1.7) {$S^{13}$};

  \draw[edge] (A) -- (B);
  \draw[edge, blue!70!black, very thick] (A) -- (C);
  \draw[edge] (A) -- (D);
  \draw[edge] (B) -- (E);
  \draw[edge] (B) -- (F);
  \draw[edge] (E) -- (J);
  \draw[edge] (F) -- (L);
  \draw[edge] (F) -- (K);
  \draw[edge] (K) -- (Q);
  \draw[edge, blue!70!black, very thick] (C) -- (H);
  \draw[edge] (C) -- (G);
  \draw[edge] (H) -- (N);
  \draw[edge, blue!70!black, very thick] (H) -- (O);
  \draw[edge, blue!70!black, very thick] (O) -- (S);
\end{tikzpicture}
\caption{DFS apos alteracao; D e G permanecem na fronteira quando S e encontrado}
\end{figure}

\textbf{Arvores IDS apos alteracao (limites 0 a 4):}
\begin{figure}[H]
\centering
\begin{subfigure}[t]{0.32\textwidth}
\centering
\begin{forest}
for tree={draw,circle,minimum size=6mm,inner sep=0pt,l sep=8mm,s sep=6mm,edge={->,thick}}
[{$A^{1}$}]
\end{forest}
\caption{Limite 0}
\end{subfigure}
\hfill
\begin{subfigure}[t]{0.32\textwidth}
\centering
\begin{forest}
for tree={draw,circle,minimum size=6mm,inner sep=0pt,l sep=8mm,s sep=6mm,edge={->,thick}}
[{$A^{1}$}
  [{$B^{2}$}]
  [{$C^{3}$}, edge+={blue!70!black,very thick}]
  [{$D^{4}$}]
]
\end{forest}
\caption{Limite 1}
\end{subfigure}
\hfill
\begin{subfigure}[t]{0.32\textwidth}
\centering
\begin{forest}
for tree={draw,circle,minimum size=6mm,inner sep=0pt,l sep=7mm,s sep=4mm,edge={->,thick}}
[{$A^{1}$}
  [{$B^{2}$}
    [{$E^{3}$}]
    [{$F^{4}$}]
  ]
  [{$C^{5}$}, edge+={blue!70!black,very thick}
    [{$H^{6}$}, edge+={blue!70!black,very thick}]
    [{$G^{7}$}]
  ]
  [{$D^{8}$}
    [{$I^{9}$}]
  ]
]
\end{forest}
\caption{Limite 2}
\end{subfigure}

\vspace{0.8em}

\begin{subfigure}[t]{0.49\textwidth}
\centering
\begin{forest}
for tree={draw,circle,minimum size=6mm,inner sep=0pt,l sep=7mm,s sep=4mm,edge={->,thick}}
[{$A^{1}$}
  [{$B^{2}$}
    [{$E^{3}$}
      [{$J^{4}$}]
    ]
    [{$F^{5}$}
      [{$L^{6}$}]
      [{$K^{7}$}]
    ]
  ]
  [{$C^{8}$}, edge+={blue!70!black,very thick}
    [{$H^{9}$}, edge+={blue!70!black,very thick}
      [{$N^{10}$}]
      [{$O^{11}$}, edge+={blue!70!black,very thick}]
    ]
    [{$G^{12}$}
      [{$M^{13}$}]
    ]
  ]
  [{$D^{14}$}
    [{$I^{15}$}
      [{$P^{16}$}]
    ]
  ]
]
\end{forest}
\caption{Limite 3}
\end{subfigure}
\hfill
\begin{subfigure}[t]{0.49\textwidth}
\centering
\begin{forest}
for tree={draw,circle,minimum size=6mm,inner sep=0pt,l sep=7mm,s sep=4mm,edge={->,thick}}
[{$A^{1}$}
  [{$B^{2}$}
    [{$E^{3}$}
      [{$J^{4}$}]
    ]
    [{$F^{5}$}
      [{$L^{6}$}]
      [{$K^{7}$}
        [{$Q^{8}$}]
      ]
    ]
  ]
  [{$C^{9}$}, edge+={blue!70!black,very thick}
    [{$H^{10}$}, edge+={blue!70!black,very thick}
      [{$N^{11}$}]
      [{$O^{12}$}, edge+={blue!70!black,very thick}
        [{$S^{13}$}, edge+={blue!70!black,very thick}]
      ]
    ]
    [G]
  ]
  [D]
]
\end{forest}
\caption{Limite 4}
\end{subfigure}
\caption{IDS apos alteracao: uma arvore por limite; caminho solucao em azul}
\end{figure}

\textbf{Mudanca na ordem de expansao (resumo):}
\begin{itemize}[leftmargin=1.5em]
  \item BFS original: $[A,B,C,D,E,F,G,H,I,J,K,L,M,N,O,P,Q,R,S]$.
  \item BFS modificado: $[A,B,C,D,E,F,H,G,I,J,L,K,N,O,M,P,Q,S]$.
  \item DFS original: $[A,B,E,J,F,K,Q,L,C,G,M,R,H,N,O,S]$.
  \item DFS modificado: $[A,B,E,J,F,L,K,Q,C,H,N,O,S]$.
  \item IDS (limite 4) original: $[A,B,E,J,F,K,Q,L,C,G,M,R,H,N,O,S]$.
  \item IDS (limite 4) modificado: $[A,B,E,J,F,L,K,Q,C,H,N,O,S]$.
\end{itemize}

\subsection{Respostas objetivas do experimento}
\begin{itemize}[leftmargin=1.5em]
  \item a) O caminho solucao mudou? \textbf{Nao}. Em todos os casos permaneceu $A \rightarrow C \rightarrow H \rightarrow O \rightarrow S$.
  \item b) A quantidade de nos expandidos mudou? \textbf{Sim}. BFS: $19 \rightarrow 18$, DFS: $16 \rightarrow 13$, IDS (total): $46 \rightarrow 43$.
  \item c) Algum algoritmo foi mais sensivel a ordem dos sucessores? \textbf{Sim, DFS}. Foi o que mais reduziu expansoes (queda de 3, cerca de 18,75\%).
  \item d) Qual algoritmo apresentou maior consumo de memoria? \textbf{BFS}. Maior fronteira maxima observada ($7$ nos), contra $4$ em DFS/IDS.
  \item e) Qual algoritmo encontrou a solucao mais rapidamente? \textbf{DFS modificado}, com 13 nos expandidos.
\end{itemize}

\subsection{Comparacao final entre os algoritmos}
\textbf{Notacao das complexidades:}
\begin{itemize}[leftmargin=1.5em]
  \item $b$: fator de ramificacao medio.
  \item $d$: profundidade da solucao mais rasa.
  \item $m$: profundidade maxima da arvore de busca.
\end{itemize}

\begin{enumerate}[label=\alph*)]
  \item \textbf{Completude}
  \begin{itemize}[leftmargin=1.5em]
    \item BFS: completo em espaco finito e com $b$ finito.
    \item DFS: nao completo em geral (pode prender em ramo muito profundo/infinito); neste problema finito, encontrou solucao.
    \item IDS: completo em espaco finito e com $b$ finito.
  \end{itemize}

  \item \textbf{Otimalidade}
  \begin{itemize}[leftmargin=1.5em]
    \item BFS: otimo quando os custos de aresta sao unitarios (caso desta questao).
    \item DFS: nao e otimo em geral.
    \item IDS: otimo com custos unitarios (encontra a primeira solucao na menor profundidade).
  \end{itemize}

  \item \textbf{Complexidade de tempo}
  \begin{itemize}[leftmargin=1.5em]
    \item BFS: $O(b^d)$.
    \item DFS: $O(b^m)$.
    \item IDS: $O(b^d)$ (com custo extra por reexpandir nos em limites menores).
  \end{itemize}

  \item \textbf{Complexidade de espaco}
  \begin{itemize}[leftmargin=1.5em]
    \item BFS: $O(b^d)$ (guarda toda a fronteira do nivel).
    \item DFS: $O(bm)$ (guarda basicamente o caminho atual e alternativas locais).
    \item IDS: $O(bd)$ (parecido com DFS por iteracao).
  \end{itemize}

  \item \textbf{Dependencia da ordem dos sucessores}
  \begin{itemize}[leftmargin=1.5em]
    \item BFS: media (a ordem afeta expansoes, mas menos que DFS).
    \item DFS: alta (a ordem muda fortemente o ramo seguido primeiro).
    \item IDS: media/alta (cada iteracao e DFS-limitada, entao herda essa sensibilidade).
  \end{itemize}

  \item \textbf{Comportamento em arvores profundas}
  \begin{itemize}[leftmargin=1.5em]
    \item BFS: tende a estourar memoria.
    \item DFS: usa pouca memoria, mas pode descer em ramos ruins.
    \item IDS: boa alternativa quando ha limite de memoria, ao custo de repeticao.
  \end{itemize}

  \item \textbf{Adequacao para problemas grandes}
  \begin{itemize}[leftmargin=1.5em]
    \item BFS: limitado quando a fronteira cresce muito.
    \item DFS: bom quando memoria e critica e otimalidade nao e obrigatoria.
    \item IDS: bom compromisso entre completude e memoria em problemas de custo unitario.
  \end{itemize}
\end{enumerate}
